\documentclass[11pt,a4paper]{article}

\usepackage[utf8]{inputenc}
\usepackage[T1]{fontenc}
\usepackage[french]{babel}
\usepackage[a4paper,margin=2.2cm]{geometry}
\usepackage{amsmath,amssymb,amsthm,mathtools}
\usepackage{lmodern}
\usepackage{enumitem}
\usepackage{hyperref}
\usepackage{xcolor}
\usepackage{fancyhdr}
\usepackage{stmaryrd}
\usepackage{mathrsfs}
\usepackage{array}
\usepackage{bm}

\hypersetup{
    colorlinks=true,
    linkcolor=black,
    urlcolor=blue,
    citecolor=black
}

\setlength{\parindent}{0pt}
\setlength{\parskip}{0.6em}

\pagestyle{fancy}
\fancyhf{}
\fancyhead[L]{Réduction --- Chapitre 18}
\fancyhead[R]{Réduction des endomorphismes symétriques}
\fancyfoot[C]{\thepage}

% Environnements
\newtheorem{definition}{Définition}[section]
\newtheorem{theoreme}[definition]{Théorème}
\newtheorem{proposition}[definition]{Proposition}
\newtheorem{corollaire}[definition]{Corollaire}
\newtheorem{lemme}[definition]{Lemme}
\theoremstyle{remark}
\newtheorem{remarque}[definition]{Remarque}
\newtheorem{exemple}[definition]{Exemple}
\newtheorem*{methode}{Méthode}
\newtheorem*{idee}{Idée}

% Commandes
\newcommand{\R}{\mathbb{R}}
\newcommand{\N}{\mathbb{N}}
\newcommand{\K}{\mathbb{K}}
\newcommand{\Vect}{\mathrm{Vect}}
\newcommand{\Ker}{\mathrm{Ker}}
\newcommand{\Img}{\mathrm{Im}}
\newcommand{\rg}{\mathrm{rg}}
\newcommand{\id}{\mathrm{Id}}
\newcommand{\Sp}{\mathrm{Sp}}
\newcommand{\tr}{\mathrm{tr}}
\newcommand{\diag}{\mathrm{diag}}
\newcommand{\Norm}[1]{\left\|#1\right\|}
\newcommand{\scal}[2]{\left\langle #1,#2\right\rangle}
\newcommand{\Mnn}{M_n(\R)}
\newcommand{\Symn}{S_n(\R)}
\newcommand{\On}{O_n(\R)}

\begin{document}

\begin{center}
    {\LARGE \textbf{Chapitre 18 --- Réduction des endomorphismes symétriques}}\\[0.4em]
    {\large Théorème spectral et diagonalisation orthogonale}
\end{center}

\hrule
\vspace{1em}

\tableofcontents

\newpage

\section{Pourquoi ce chapitre est-il un aboutissement ?}

Dans les chapitres précédents, on a étudié :
\begin{itemize}
    \item les valeurs propres ;
    \item la diagonalisation ;
    \item la trigonalisation ;
    \item les espaces euclidiens.
\end{itemize}

Ici, ces thèmes se rejoignent de façon particulièrement belle.

Les matrices symétriques réelles et les endomorphismes symétriques ont des propriétés remarquables :
\begin{itemize}
    \item leurs valeurs propres sont réelles ;
    \item leurs sous-espaces propres associés à des valeurs propres distinctes sont orthogonaux ;
    \item ils sont toujours diagonalisables ;
    \item mieux encore : ils sont \textbf{orthogonalement diagonalisables}.
\end{itemize}

\begin{idee}
Le cas symétrique est le cas où l’algèbre linéaire et la géométrie euclidienne s’accordent parfaitement.
\end{idee}

Ce chapitre mène au \textbf{théorème spectral}, l’un des résultats majeurs de l’algèbre linéaire réelle.

\section{Rappels sur les espaces euclidiens}

Soit \(E\) un espace euclidien réel, c’est-à-dire un espace vectoriel réel de dimension finie muni d’un produit scalaire
\[
\scal{\cdot}{\cdot}:E\times E\to \R.
\]

On rappelle que :
\begin{itemize}
    \item une base orthonormée est une base \((e_1,\dots,e_n)\) telle que
    \[
    \scal{e_i}{e_j}=\delta_{ij} ;
    \]
    \item un endomorphisme \(u\) est représenté dans une base orthonormée par une matrice réelle ;
    \item une matrice \(P\in M_n(\R)\) est orthogonale si
    \[
    P^TP=I_n.
    \]
\end{itemize}

\section{Endomorphismes symétriques}

\subsection{Définition}

\begin{definition}
Soit \(E\) un espace euclidien réel. Un endomorphisme
\[
u\in\mathcal L(E,E)
\]
est dit \textbf{symétrique} si
\[
\forall x,y\in E,\qquad \scal{u(x)}{y}=\scal{x}{u(y)}.
\]
\end{definition}

\begin{remarque}
Cette définition signifie que \(u\) est auto-adjoint pour le produit scalaire.
\end{remarque}

\subsection{Lien matriciel}

\begin{proposition}
Dans une base orthonormée, un endomorphisme est symétrique si et seulement si sa matrice est symétrique, c’est-à-dire vérifie
\[
A^T=A.
\]
\end{proposition}

\begin{proof}
Soit \((e_1,\dots,e_n)\) une base orthonormée de \(E\), et \(A=(a_{ij})\) la matrice de \(u\) dans cette base.

On a
\[
a_{ij}=\scal{u(e_j)}{e_i}.
\]
Donc \(u\) est symétrique si et seulement si, pour tous \(i,j\),
\[
\scal{u(e_j)}{e_i}=\scal{e_j}{u(e_i)}.
\]
Or
\[
\scal{e_j}{u(e_i)}=\scal{u(e_i)}{e_j}=a_{ji},
\]
d’où
\[
a_{ij}=a_{ji}.
\]
Cela équivaut à
\[
A^T=A.
\]
\end{proof}

\begin{corollaire}
La notion de matrice symétrique réelle est la traduction matricielle de la notion d’endomorphisme symétrique dans une base orthonormée.
\end{corollaire}

\section{Premières propriétés}

\subsection{Stabilité de l’orthogonal d’un sous-espace stable}

\begin{proposition}
Soit \(u\) un endomorphisme symétrique de \(E\), et soit \(F\) un sous-espace stable par \(u\). Alors \(F^\Orth\) est aussi stable par \(u\).
\end{proposition}

\begin{proof}
Soit \(x\in F^\Orth\). Il faut montrer que \(u(x)\in F^\Orth\), c’est-à-dire que pour tout \(y\in F\),
\[
\scal{u(x)}{y}=0.
\]
Or, par symétrie de \(u\),
\[
\scal{u(x)}{y}=\scal{x}{u(y)}.
\]
Comme \(F\) est stable, \(u(y)\in F\). Et comme \(x\in F^\Orth\), on a
\[
\scal{x}{u(y)}=0.
\]
Donc
\[
\scal{u(x)}{y}=0
\]
pour tout \(y\in F\). Ainsi \(u(x)\in F^\Orth\).
\end{proof}

\begin{remarque}
Cette propriété est très importante dans la démonstration du théorème spectral par récurrence.
\end{remarque}

\subsection{Orthogonalité des sous-espaces propres}

\begin{theoreme}
Soit \(u\) un endomorphisme symétrique. Si \(\lambda\) et \(\mu\) sont deux valeurs propres distinctes de \(u\), alors les sous-espaces propres \(E_\lambda\) et \(E_\mu\) sont orthogonaux.
\end{theoreme}

\begin{proof}
Soient \(x\in E_\lambda\) et \(y\in E_\mu\). Alors
\[
u(x)=\lambda x
\qquad \text{et} \qquad
u(y)=\mu y.
\]
Comme \(u\) est symétrique,
\[
\scal{u(x)}{y}=\scal{x}{u(y)}.
\]
Donc
\[
\lambda \scal{x}{y}=\mu \scal{x}{y}.
\]
Ainsi
\[
(\lambda-\mu)\scal{x}{y}=0.
\]
Comme \(\lambda\neq \mu\), on obtient
\[
\scal{x}{y}=0.
\]
Donc \(E_\lambda\Orth E_\mu\).
\end{proof}

\begin{corollaire}
Des vecteurs propres associés à des valeurs propres distinctes d’un endomorphisme symétrique sont orthogonaux.
\end{corollaire}

\begin{remarque}
C’est beaucoup plus fort que la simple indépendance linéaire obtenue dans le cas général.
\end{remarque}

\section{Existence de valeurs propres réelles}

\subsection{Cas matriciel réel symétrique}

\begin{theoreme}
Toute matrice symétrique réelle possède au moins une valeur propre réelle.
\end{theoreme}

\begin{proof}
Soit \(A\in S_n(\R)\). Son polynôme caractéristique \(\chi_A\) est un polynôme réel de degré \(n\). Considéré sur \(\C\), il admet au moins une racine complexe.

Mais pour une matrice symétrique réelle, toute valeur propre complexe est en fait réelle. En effet, si
\[
AX=\lambda X
\]
avec \(X\neq 0\) dans \(\C^n\), alors en utilisant le produit hermitien canonique,
\[
\langle AX,X\rangle = \lambda \langle X,X\rangle.
\]
Or, comme \(A\) est réelle symétrique, on a aussi
\[
\langle AX,X\rangle = \langle X,AX\rangle = \overline{\lambda}\langle X,X\rangle.
\]
Donc
\[
\lambda \langle X,X\rangle = \overline{\lambda}\langle X,X\rangle.
\]
Comme \(\langle X,X\rangle>0\), on obtient
\[
\lambda=\overline{\lambda},
\]
donc \(\lambda\in\R\).
\end{proof}

\begin{remarque}
Dans le cadre d’un cours purement réel, on peut admettre ce résultat comme point de départ de la démonstration du théorème spectral.
\end{remarque}

\subsection{Version endomorphisme}

\begin{corollaire}
Tout endomorphisme symétrique d’un espace euclidien non nul admet au moins une valeur propre réelle.
\end{corollaire}

\begin{proof}
On choisit une base orthonormée, dans laquelle sa matrice est réelle symétrique. On applique le théorème précédent.
\end{proof}

\section{Théorème spectral}

\subsection{Énoncé}

\begin{theoreme}[Théorème spectral]
Soit \(E\) un espace euclidien réel de dimension finie, et \(u\in \mathcal L(E,E)\) un endomorphisme symétrique. Alors il existe une base orthonormée de \(E\) formée de vecteurs propres de \(u\).

Autrement dit, tout endomorphisme symétrique est diagonalisable dans une base orthonormée.
\end{theoreme}

\subsection{Démonstration}

\begin{proof}
On raisonne par récurrence sur \(n=\dim(E)\).

\textbf{Initialisation :} si \(n=1\), le résultat est évident.

\textbf{Hérédité :} supposons le résultat vrai jusqu’à la dimension \(n-1\), et soit \(E\) de dimension \(n\).

Comme \(u\) est symétrique, il admet une valeur propre réelle \(\lambda\). Soit \(e_1\) un vecteur propre associé, choisi de norme \(1\). Alors
\[
u(e_1)=\lambda e_1.
\]
Le sous-espace
\[
F=\Vect(e_1)
\]
est stable par \(u\). Par la proposition précédente, son orthogonal \(F^\Orth\) est aussi stable par \(u\).

Or \(F^\Orth\) est de dimension \(n-1\), et la restriction
\[
u_{|F^\Orth}
\]
est encore symétrique. Par hypothèse de récurrence, il existe une base orthonormée
\[
(e_2,\dots,e_n)
\]
de \(F^\Orth\) constituée de vecteurs propres de \(u_{|F^\Orth}\).

Alors
\[
(e_1,e_2,\dots,e_n)
\]
est une base orthonormée de \(E\), et chaque vecteur est propre pour \(u\). Le résultat est démontré.
\end{proof}

\begin{corollaire}
Tout endomorphisme symétrique est diagonalisable.
\end{corollaire}

\begin{proof}
Une base orthonormée est en particulier une base.
\end{proof}

\begin{corollaire}
Toute matrice symétrique réelle est diagonalisable dans une base orthonormée de \(\R^n\).
\end{corollaire}

\section{Diagonalisation orthogonale}

\subsection{Énoncé matriciel}

\begin{theoreme}
Une matrice réelle \(A\in M_n(\R)\) est symétrique si et seulement s’il existe une matrice orthogonale \(P\in O_n(\R)\) et une matrice diagonale réelle \(D\) telles que
\[
A=PDP^T.
\]
\end{theoreme}

\begin{proof}
Supposons d’abord \(A\) symétrique. Par le théorème spectral, il existe une base orthonormée de \(\R^n\) formée de vecteurs propres de \(A\). Si l’on place ces vecteurs en colonnes dans une matrice \(P\), alors \(P\) est orthogonale, et dans cette base la matrice de \(A\) est diagonale \(D\). On a donc
\[
A=PDP^{-1}=PDP^T.
\]

Réciproquement, si
\[
A=PDP^T
\]
avec \(P\) orthogonale et \(D\) diagonale réelle, alors
\[
A^T=(PDP^T)^T=P D^T P^T=PDP^T=A,
\]
car \(D^T=D\).
\end{proof}

\begin{definition}
Une telle écriture est appelée une \textbf{diagonalisation orthogonale} de \(A\).
\end{definition}

\begin{remarque}
La diagonalisation orthogonale est bien plus forte qu’une simple diagonalisation :
\begin{itemize}
    \item elle donne une base orthonormée ;
    \item la matrice de passage vérifie
    \[
    P^{-1}=P^T.
    \]
\end{itemize}
\end{remarque}

\section{Conséquences du théorème spectral}

\subsection{Décomposition orthogonale}

\begin{proposition}
Si \(u\) est symétrique et si \(\lambda_1,\dots,\lambda_r\) sont ses valeurs propres distinctes, alors
\[
E=E_{\lambda_1}\Orth \cdots \Orth E_{\lambda_r}.
\]
\end{proposition}

\begin{proof}
Les sous-espaces propres associés à des valeurs propres distinctes sont orthogonaux, et le théorème spectral donne une base orthonormée de vecteurs propres, donc leur somme vaut \(E\).
\end{proof}

\subsection{Trace et déterminant}

\begin{proposition}
Si \(A\) est symétrique réelle, de valeurs propres \(\lambda_1,\dots,\lambda_n\) comptées avec multiplicité, alors
\[
\tr(A)=\lambda_1+\cdots+\lambda_n
\]
et
\[
\det(A)=\lambda_1\cdots \lambda_n.
\]
\end{proposition}

\begin{proof}
Comme \(A\) est orthogonalement diagonalisable,
\[
A=PDP^T
\]
avec
\[
D=\diag(\lambda_1,\dots,\lambda_n).
\]
La trace et le déterminant sont invariants par similitude, donc
\[
\tr(A)=\tr(D)=\lambda_1+\cdots+\lambda_n,
\]
\[
\det(A)=\det(D)=\lambda_1\cdots \lambda_n.
\]
\end{proof}

\subsection{Signe des valeurs propres}

\begin{proposition}
Si \(u\) est symétrique positive, c’est-à-dire si
\[
\forall x\in E,\qquad \scal{u(x)}{x}\ge 0,
\]
alors toutes ses valeurs propres sont positives ou nulles.
\end{proposition}

\begin{proof}
Soit \(\lambda\) une valeur propre, et \(x\neq 0\) un vecteur propre associé. Alors
\[
\scal{u(x)}{x}=\scal{\lambda x}{x}=\lambda \scal{x}{x}=\lambda \Norm{x}^2.
\]
Comme \(\Norm{x}^2>0\) et \(\scal{u(x)}{x}\ge 0\), on obtient
\[
\lambda\ge 0.
\]
\end{proof}

\begin{corollaire}
Si \(u\) est symétrique définie positive, c’est-à-dire si
\[
\forall x\neq 0,\qquad \scal{u(x)}{x}>0,
\]
alors toutes ses valeurs propres sont strictement positives.
\end{corollaire}

\section{Formes quadratiques associées}

\subsection{Définition}

\begin{definition}
À toute matrice symétrique réelle \(A\in S_n(\R)\), on associe la forme quadratique
\[
q(X)=X^TAX.
\]
\end{definition}

\begin{remarque}
La symétrie de \(A\) est le cadre naturel des formes quadratiques réelles.
\end{remarque}

\subsection{Réduction orthogonale}

\begin{theoreme}
Soit \(A\in S_n(\R)\). Il existe une matrice orthogonale \(P\) telle que
\[
P^TAP=D
\]
soit diagonale. Ainsi, dans le nouveau repère orthonormé, la forme quadratique \(q\) s’écrit
\[
q(Y)=\lambda_1 y_1^2+\cdots+\lambda_n y_n^2.
\]
\end{theoreme}

\begin{proof}
Par le théorème spectral,
\[
A=PDP^T
\]
avec \(P\) orthogonale. Alors
\[
q(X)=X^TAX=X^T P D P^T X.
\]
En posant
\[
Y=P^TX,
\]
on obtient
\[
q(X)=Y^TDY=\lambda_1 y_1^2+\cdots+\lambda_n y_n^2.
\]
\end{proof}

\begin{remarque}
Cette réduction est fondamentale en géométrie et dans l’étude locale des fonctions de plusieurs variables.
\end{remarque}

\section{Interprétation géométrique}

\subsection{Axes propres}

\begin{remarque}
Les vecteurs propres d’un endomorphisme symétrique définissent des directions orthogonales privilégiées.  
Dans une base orthonormée de vecteurs propres, l’endomorphisme agit séparément sur chaque axe :
\[
u(e_i)=\lambda_i e_i.
\]
\end{remarque}

\subsection{Quadriques}

\begin{remarque}
Les équations quadratiques de la forme
\[
X^TAX=1
\]
avec \(A\) symétrique peuvent être simplifiées par diagonalisation orthogonale. On obtient alors une équation de la forme
\[
\lambda_1 y_1^2+\cdots+\lambda_n y_n^2=1,
\]
qui permet d’identifier la nature géométrique de la surface ou de la courbe.
\end{remarque}

\section{Exemples détaillés}

\subsection{Exemple simple}

\begin{exemple}
Considérons
\[
A=
\begin{pmatrix}
2 & 0\\
0 & 3
\end{pmatrix}.
\]
La matrice est déjà diagonale et symétrique. Une base orthonormée de vecteurs propres est la base canonique, et les valeurs propres sont \(2\) et \(3\).
\end{exemple}

\subsection{Exemple non diagonale}

\begin{exemple}
Considérons
\[
A=
\begin{pmatrix}
1 & 1\\
1 & 1
\end{pmatrix}.
\]
On calcule
\[
\chi_A(X)=\det\begin{pmatrix}
X-1 & -1\\
-1 & X-1
\end{pmatrix}
=(X-1)^2-1=X(X-2).
\]
Les valeurs propres sont \(0\) et \(2\).

Pour \(\lambda=2\),
\[
(A-2I_2)X=0
\Longrightarrow
\begin{pmatrix}
-1 & 1\\
1 & -1
\end{pmatrix}
X=0,
\]
donc
\[
E_2=\Vect\left(
\begin{pmatrix}
1\\1
\end{pmatrix}
\right).
\]

Pour \(\lambda=0\),
\[
AX=0
\Longrightarrow
\begin{pmatrix}
1 & 1\\
1 & 1
\end{pmatrix}
X=0,
\]
donc
\[
E_0=\Vect\left(
\begin{pmatrix}
1\\-1
\end{pmatrix}
\right).
\]

Ces deux directions propres sont orthogonales. En normalisant, on obtient une base orthonormée
\[
\left(
\frac{1}{\sqrt2}\begin{pmatrix}1\\1\end{pmatrix},
\frac{1}{\sqrt2}\begin{pmatrix}1\\-1\end{pmatrix}
\right),
\]
dans laquelle la matrice devient
\[
\diag(2,0).
\]
\end{exemple}

\subsection{Exemple de forme quadratique}

\begin{exemple}
Considérons la forme quadratique
\[
q(x,y)=x^2+2xy+y^2.
\]
La matrice symétrique associée est
\[
A=
\begin{pmatrix}
1 & 1\\
1 & 1
\end{pmatrix}.
\]
Comme dans l’exemple précédent, elle se diagonalise orthogonalement en
\[
\diag(2,0).
\]
Dans une base orthonormée adaptée,
\[
q(y_1,y_2)=2y_1^2.
\]
\end{exemple}

\section{Méthode pratique de diagonalisation orthogonale}

\begin{methode}
Pour orthogonalement diagonaliser une matrice symétrique réelle \(A\) :
\begin{enumerate}[label=\arabic*)]
    \item vérifier que
    \[
    A^T=A ;
    \]
    \item calculer le polynôme caractéristique ;
    \item déterminer les valeurs propres ;
    \item calculer les sous-espaces propres ;
    \item choisir dans chaque sous-espace propre une base orthonormée ;
    \item réunir ces bases : on obtient une base orthonormée de \(\R^n\) ;
    \item former la matrice orthogonale \(P\) ayant ces vecteurs propres orthonormés pour colonnes ;
    \item écrire
    \[
    A=PDP^T.
    \]
\end{enumerate}
\end{methode}

\begin{remarque}
Les sous-espaces propres associés à des valeurs propres distinctes étant orthogonaux, il suffit d’orthonormaliser séparément à l’intérieur de chacun d’eux.
\end{remarque}

\section{Matrices orthogonales et changement de base}

\begin{proposition}
Soit \(P\in O_n(\R)\). Alors le changement de base associé conserve le produit scalaire, les normes et les angles.
\end{proposition}

\begin{proof}
Pour tous \(X,Y\in\R^n\),
\[
(PX)^T(PY)=X^T(P^TP)Y=X^TY.
\]
Donc le produit scalaire est conservé. Les normes et angles aussi.
\end{proof}

\begin{remarque}
La diagonalisation orthogonale n’est pas seulement une simplification algébrique : elle respecte entièrement la géométrie euclidienne.
\end{remarque}

\section{Critères de positivité}

\subsection{Via les valeurs propres}

\begin{theoreme}
Soit \(A\in S_n(\R)\). Alors :
\begin{enumerate}[label=\arabic*)]
    \item \(A\) est positive si et seulement si toutes ses valeurs propres sont positives ou nulles ;
    \item \(A\) est définie positive si et seulement si toutes ses valeurs propres sont strictement positives.
\end{enumerate}
\end{theoreme}

\begin{proof}
La première implication a déjà été démontrée.

Réciproquement, si \(A=PDP^T\) avec \(D=\diag(\lambda_1,\dots,\lambda_n)\) et \(\lambda_i\ge 0\), alors pour tout \(X\),
\[
X^TAX=X^TPDP^TX=Y^TDY=\sum_{i=1}^n \lambda_i y_i^2\ge 0,
\]
où \(Y=P^TX\). Donc \(A\) est positive.

Le cas définie positive est analogue.
\end{proof}

\section{Erreurs classiques}

\begin{itemize}
    \item Confondre matrice symétrique et matrice orthogonale.
    \item Oublier que le théorème spectral est un résultat réel : il porte sur les matrices symétriques réelles et les espaces euclidiens réels.
    \item Diagonaliser une matrice symétrique sans imposer que la base de vecteurs propres soit orthonormée, alors que c’est justement l’intérêt du théorème.
    \item Oublier de normaliser les vecteurs propres avant de construire \(P\).
    \item Se tromper entre
    \[
    A=PDP^{-1}
    \quad \text{et} \quad
    A=PDP^T.
    \]
\end{itemize}

\section{Résumé du chapitre}

\begin{itemize}
    \item Un endomorphisme symétrique vérifie
    \[
    \scal{u(x)}{y}=\scal{x}{u(y)}.
    \]
    \item Dans une base orthonormée, sa matrice est symétrique :
    \[
    A^T=A.
    \]
    \item Les sous-espaces propres associés à des valeurs propres distinctes sont orthogonaux.
    \item Toute matrice symétrique réelle admet des valeurs propres réelles.
    \item Théorème spectral : tout endomorphisme symétrique réel admet une base orthonormée de vecteurs propres.
    \item Toute matrice symétrique réelle se diagonalise orthogonalement :
    \[
    A=PDP^T.
    \]
    \item Les formes quadratiques associées aux matrices symétriques se réduisent en somme de carrés dans une base orthonormée adaptée.
\end{itemize}

\section*{Exercices d’application immédiate}

\begin{enumerate}[label=\textbf{Exercice \arabic*.}]
    \item Montrer que la matrice
    \[
    \begin{pmatrix}
    2 & 1\\
    1 & 2
    \end{pmatrix}
    \]
    est symétrique.

    \item Calculer les valeurs propres de
    \[
    \begin{pmatrix}
    1 & 1\\
    1 & 1
    \end{pmatrix}.
    \]

    \item Déterminer une base orthonormée de vecteurs propres de la matrice précédente.

    \item Orthogonalement diagonaliser
    \[
    \begin{pmatrix}
    2 & 0\\
    0 & 3
    \end{pmatrix}.
    \]

    \item Montrer que les sous-espaces propres associés à des valeurs propres distinctes d’un endomorphisme symétrique sont orthogonaux.

    \item Montrer qu’une matrice symétrique réelle positive a toutes ses valeurs propres positives ou nulles.

    \item Réduire la forme quadratique
    \[
    q(x,y)=x^2+2xy+y^2
    \]
    par diagonalisation orthogonale.

    \item Montrer qu’une matrice orthogonalement diagonalisable est nécessairement symétrique.

    \item Soit
    \[
    A=
    \begin{pmatrix}
    3 & 1\\
    1 & 3
    \end{pmatrix}.
    \]
    Déterminer ses valeurs propres et l’orthogonalement diagonaliser.

    \item Montrer qu’une matrice symétrique réelle est définie positive si et seulement si toutes ses valeurs propres sont strictement positives.
\end{enumerate}

\end{document}